% ============================================================
% 2. RELATED WORK
% ============================================================
\section{Related Work}
\label{sec:related}

\paragraph{Adaptive test-time compute.}
Adaptive gating methods reduce the overhead of test-time
optimization by invoking it only when needed. In the
\emph{reasoning} setting, signal-based methods use token-level
confidence or entropy to decide when to extend
thinking~\citep{seag2025, cats2025, corefine2026,
thinkjustenough2025, s1_2025, han2025tokenbudget}. RL-based methods
learn think/no-think policies through reinforcement
learning~\citep{adaptthink2025, thinkless2025, aggarwal2025l1},
though these require per-environment training and provide no
interpretable direction signal. Hidden-state probing
methods~\citep{diffadapt2025, zhu2025llmalreadyknows} train
lightweight classifiers on internal representations to estimate task
difficulty, assuming a universal difficulty-to-compute mapping.
In the \emph{agent} setting, vote-based
methods~\citep{catts2026} sample multiple candidate actions and
trigger evaluation when disagreement is high, and
\citet{paglieri2025learning} learn when to invoke a planner at the
episode level. The vast majority of adaptive compute research
operates on homogeneous reasoning benchmarks (MATH, GSM8K, GPQA),
where signal direction happens to be consistent across tasks. Our
work is the first to systematically study signal semantics across
heterogeneous \emph{agent} environments, where we find that the
direction reverses. A detailed per-method comparison is provided in
Appendix~\ref{app:related}.

\paragraph{Test-time scaling and search.}
Test-time scaling~\citep{snell2024scaling} studies compute-optimal
allocation at the problem level, showing that adaptive allocation
outperforms uniform allocation by up to 4$\times$ in efficiency.
This line of work focuses on \emph{how much} compute to allocate
given a fixed difficulty signal, whereas we question the signal
itself. Search methods such as Tree of
Thoughts~\citep{yao2023tree}, LATS~\citep{zhou2024lats},
RAP~\citep{hao2023reasoning}, and FLARE~\citep{wang2026flare}
improve the quality of the search process through tree search,
value propagation, and lookahead evaluation. These are complementary
to our work: they improve \emph{what} the optimizer does, while we
control \emph{when} to invoke it. Our gating framework can be
composed with any such optimizer.

\paragraph{Uncertainty estimation and metareasoning.}
Recent work on LLM uncertainty~\citep{tao2025revisiting,
heo2025llmuncertainty} observes that calibration quality varies
substantially across task types, providing converging evidence that
uncertainty semantics are domain-dependent. \citet{yadkori2024believe}
propose decoupling epistemic and aleatoric uncertainty in LLM
outputs, a distinction conceptually related to our Type~I/Type~D
decomposition. From a decision-theoretic perspective, rational
metareasoning~\citep{russell1991right, desabbata2024metareasoning}
formalizes the value of computation (VOC) as a guide for allocating
cognitive resources. Our work extends this framework by showing
that VOC can be \emph{negative} in the agent setting when the
evaluator and executor are the same model, and that direction
discovery is a necessary condition for non-negative VOC across
environments.
